% grow=2.5 is not a universal area term. LOO kills it.
\documentclass[11pt]{article}
\usepackage[margin=1.05in]{geometry}
\usepackage{amsmath,amssymb,amsthm}
\usepackage[colorlinks=true,linkcolor=blue,citecolor=blue,urlcolor=blue]{hyperref}

\newtheorem*{admission}{Admission}

\title{\textbf{The Leftover $2.5$ Is Not an Area Law.\\
Held-Out Rods Miss by Five.}}
\author{Robert W.\ McGwier\\
\small CTO, Cohere Technology Group\\
\small GitHub: \texttt{n4hy}}
\date{15 August 2026 --- Version 1}

\begin{document}
\maketitle

\begin{abstract}
\noindent
Paper~64 left a residual: slab grow is $2.5$, not $3$.
This note asks whether
\[
\delta S=aV+bA
\]
is a universal area term, or shape junk.

$N=12$. Balls, cubes, slabs, rods.
$\alpha=0.01,0.02,0.04$.
In-sample two-term rel RMS $0.081,0.073,0.066$.
Volume-only is the same to two digits.
Leave-one-family-out: held-out rods miss by
$5.2$. The area coefficient flips sign when
slabs are withheld. Auditor $N=10$: gain and
LOO refuted.

There is no Clausius/horizon handle in this
residual. Sea-transfer $\delta S$ is volume-like
with $O(1)$ shape junk. Gravity remains
$\sum\delta e$ plus inherited Gauss.
\end{abstract}

\begin{admission}
I hoped grow $\neq 3$ was an area term.
Leave-one-family-out killed it. I did not
move C\_loo.
\end{admission}

\end{document}
